%% ---------------------------------------------------------------------------
\section{Discussion}\label{sec:discussion}

The two barriers constrain the half-line from opposite sides, and the
five-state rules of \cref{sec:five} show how they complement each other. The
half-line of $\delta_{12}$ meets the speed-$1/3$ barrier exactly in simulations,
but its periodic region at the left border is excluded by the left-border
barrier (\cref{cor:germ12}). From cell $3$ on, the half-line of $\delta_{14}$
evolves under the elementary cellular automaton $22$ on the states $\Lq$ and
$\Bq$; it passes both barriers in the simulated ranges, and the rule fails only
in the first rows of the reflected triangle of the line of length $15$. All five rules of
\cref{prop:nearmiss} fail in the first rows of their reflected triangles.

The early rows and the right ends of the reflected triangles are controlled by
\cref{cor:reduce,lem:G}, but they do not exclude every rule. For the germ of
$\delta_{14}^{\mathrm d}$, which has complexity $15$, conditions (G1)--(G3) with
$T=10$ and all $n\ge12$ have a model, and so do (G1)--(G3) with $T=7$ and all
$n\ge9$ together with (G5) for $9\le n\le20$. In the second model, the lines of
lengths $9$, $11,\dots,17$ and $19$ have both end cells, but not all cells, in
state $\Fq$ at time $2n-2$, and no cell in state $\Fq$ before
(\cref{fig:model}(a); both models are in the supplementary material). This suggests that excluding such rules needs
conditions on the middle cells of the firing row, which are expensive to
encode: for $n=15$, the dependency cone of the middle cell of the firing row
contains $76$ of the $120$ cells of $R_{15}$ (\cref{fig:model}(b)).

\begin{figure}[!htb]
\centering
%% Figure: the second model of Section 7 on the line of length 15, and two dependency cones of R_15
%% (cells and cones from figures.py). In (b), the relative cell (tau,kappa) is [14-kappa,15-kappa] x [-tau-1,-tau],
%% shifted so that it is level with the cell (14+tau, 15-kappa) of (a).
\begin{tikzpicture}[x=1.55mm,y=1.55mm]
  \def\styleswatch#1{\tikz[baseline=-0.6ex,x=1cm,y=1cm]\filldraw[#1,draw=grid,line width=0.3pt] (0,-0.09) rectangle (0.24,0.15);}
  % (a) the line of length 15 up to time 28
  \begin{scope}
    %% generated by figures.py; do not edit
\fill[stG] (0,0) rectangle (1,-1);
\fill[stA] (0,-1) rectangle (1,-2);
\fill[stB] (1,-1) rectangle (2,-2);
\fill[stA] (0,-2) rectangle (1,-3);
\fill[stB] (1,-2) rectangle (3,-3);
\fill[stA] (0,-3) rectangle (1,-4);
\fill[stB] (1,-3) rectangle (2,-4);
\fill[stB] (3,-3) rectangle (4,-4);
\fill[stA] (0,-4) rectangle (1,-5);
\fill[stB] (1,-4) rectangle (3,-5);
\fill[stA] (3,-4) rectangle (4,-5);
\fill[stB] (4,-4) rectangle (5,-5);
\fill[stA] (0,-5) rectangle (1,-6);
\fill[stB] (1,-5) rectangle (2,-6);
\fill[stB] (3,-5) rectangle (6,-6);
\fill[stA] (0,-6) rectangle (1,-7);
\fill[stB] (1,-6) rectangle (3,-7);
\fill[stB] (4,-6) rectangle (5,-7);
\fill[stB] (6,-6) rectangle (7,-7);
\fill[stA] (0,-7) rectangle (1,-8);
\fill[stB] (1,-7) rectangle (2,-8);
\fill[stB] (3,-7) rectangle (4,-8);
\fill[stA] (4,-7) rectangle (5,-8);
\fill[stB] (5,-7) rectangle (6,-8);
\fill[stA] (6,-7) rectangle (7,-8);
\fill[stB] (7,-7) rectangle (8,-8);
\fill[stA] (0,-8) rectangle (1,-9);
\fill[stB] (1,-8) rectangle (5,-9);
\fill[stA] (5,-8) rectangle (6,-9);
\fill[stB] (6,-8) rectangle (9,-9);
\fill[stA] (0,-9) rectangle (1,-10);
\fill[stB] (1,-9) rectangle (4,-10);
\fill[stB] (5,-9) rectangle (8,-10);
\fill[stB] (9,-9) rectangle (10,-10);
\fill[stA] (0,-10) rectangle (1,-11);
\fill[stB] (1,-10) rectangle (3,-11);
\fill[stB] (4,-10) rectangle (5,-11);
\fill[stB] (6,-10) rectangle (7,-11);
\fill[stB] (8,-10) rectangle (9,-11);
\fill[stA] (9,-10) rectangle (10,-11);
\fill[stB] (10,-10) rectangle (11,-11);
\fill[stA] (0,-11) rectangle (1,-12);
\fill[stB] (1,-11) rectangle (2,-12);
\fill[stB] (3,-11) rectangle (4,-12);
\fill[stA] (4,-11) rectangle (5,-12);
\fill[stB] (5,-11) rectangle (6,-12);
\fill[stA] (6,-11) rectangle (7,-12);
\fill[stB] (7,-11) rectangle (12,-12);
\fill[stA] (0,-12) rectangle (1,-13);
\fill[stB] (1,-12) rectangle (5,-13);
\fill[stA] (5,-12) rectangle (6,-13);
\fill[stB] (6,-12) rectangle (11,-13);
\fill[stB] (12,-12) rectangle (13,-13);
\fill[stA] (0,-13) rectangle (1,-14);
\fill[stB] (1,-13) rectangle (4,-14);
\fill[stB] (5,-13) rectangle (10,-14);
\fill[stB] (11,-13) rectangle (12,-14);
\fill[stA] (12,-13) rectangle (13,-14);
\fill[stB] (13,-13) rectangle (14,-14);
\fill[stA] (0,-14) rectangle (1,-15);
\fill[stB] (1,-14) rectangle (3,-15);
\fill[stB] (4,-14) rectangle (5,-15);
\fill[stB] (6,-14) rectangle (9,-15);
\fill[stB] (10,-14) rectangle (14,-15);
\fill[stG] (14,-14) rectangle (15,-15);
\fill[stA] (0,-15) rectangle (1,-16);
\fill[stB] (1,-15) rectangle (2,-16);
\fill[stB] (3,-15) rectangle (4,-16);
\fill[stA] (4,-15) rectangle (5,-16);
\fill[stB] (5,-15) rectangle (6,-16);
\fill[stB] (7,-15) rectangle (8,-16);
\fill[stB] (9,-15) rectangle (10,-16);
\fill[stB] (11,-15) rectangle (13,-16);
\fill[stG] (13,-15) rectangle (14,-16);
\fill[stA] (14,-15) rectangle (15,-16);
\fill[stA] (0,-16) rectangle (1,-17);
\fill[stB] (1,-16) rectangle (7,-17);
\fill[stA] (7,-16) rectangle (8,-17);
\fill[stB] (8,-16) rectangle (9,-17);
\fill[stA] (9,-16) rectangle (10,-17);
\fill[stB] (10,-16) rectangle (11,-17);
\fill[stG] (12,-16) rectangle (13,-17);
\fill[stA] (13,-16) rectangle (14,-17);
\fill[stG] (14,-16) rectangle (15,-17);
\fill[stA] (0,-17) rectangle (1,-18);
\fill[stB] (1,-17) rectangle (6,-18);
\fill[stB] (7,-17) rectangle (8,-18);
\fill[stA] (8,-17) rectangle (9,-18);
\fill[stB] (9,-17) rectangle (11,-18);
\fill[stG] (11,-17) rectangle (12,-18);
\fill[stB] (12,-17) rectangle (13,-18);
\fill[stG] (13,-17) rectangle (15,-18);
\fill[stA] (0,-18) rectangle (1,-19);
\fill[stB] (1,-18) rectangle (5,-19);
\fill[stB] (6,-18) rectangle (10,-19);
\fill[stG] (10,-18) rectangle (11,-19);
\fill[stB] (11,-18) rectangle (13,-19);
\fill[stB] (14,-18) rectangle (15,-19);
\fill[stA] (0,-19) rectangle (1,-20);
\fill[stB] (1,-19) rectangle (4,-20);
\fill[stB] (5,-19) rectangle (6,-20);
\fill[stB] (7,-19) rectangle (9,-20);
\fill[stG] (9,-19) rectangle (10,-20);
\fill[stB] (10,-19) rectangle (11,-20);
\fill[stA] (11,-19) rectangle (12,-20);
\fill[stB] (13,-19) rectangle (15,-20);
\fill[stA] (0,-20) rectangle (1,-21);
\fill[stB] (1,-20) rectangle (3,-21);
\fill[stB] (4,-20) rectangle (5,-21);
\fill[stA] (5,-20) rectangle (6,-21);
\fill[stB] (6,-20) rectangle (7,-21);
\fill[stG] (8,-20) rectangle (9,-21);
\fill[stB] (9,-20) rectangle (10,-21);
\fill[stG] (10,-20) rectangle (11,-21);
\fill[stB] (11,-20) rectangle (12,-21);
\fill[stG] (12,-20) rectangle (13,-21);
\fill[stA] (14,-20) rectangle (15,-21);
\fill[stA] (0,-21) rectangle (1,-22);
\fill[stB] (1,-21) rectangle (2,-22);
\fill[stB] (3,-21) rectangle (7,-22);
\fill[stG] (7,-21) rectangle (8,-22);
\fill[stA] (8,-21) rectangle (9,-22);
\fill[stB] (9,-21) rectangle (12,-22);
\fill[stA] (12,-21) rectangle (14,-22);
\fill[stA] (0,-22) rectangle (1,-23);
\fill[stB] (1,-22) rectangle (3,-23);
\fill[stB] (4,-22) rectangle (6,-23);
\fill[stG] (6,-22) rectangle (7,-23);
\fill[stA] (7,-22) rectangle (9,-23);
\fill[stB] (9,-22) rectangle (11,-23);
\fill[stB] (13,-22) rectangle (14,-23);
\fill[stA] (0,-23) rectangle (1,-24);
\fill[stB] (1,-23) rectangle (2,-24);
\fill[stB] (3,-23) rectangle (4,-24);
\fill[stG] (5,-23) rectangle (6,-24);
\fill[stA] (6,-23) rectangle (7,-24);
\fill[stB] (7,-23) rectangle (8,-24);
\fill[stB] (9,-23) rectangle (10,-24);
\fill[stB] (11,-23) rectangle (12,-24);
\fill[stG] (12,-23) rectangle (13,-24);
\fill[stA] (13,-23) rectangle (14,-24);
\fill[stG] (14,-23) rectangle (15,-24);
\fill[stA] (0,-24) rectangle (1,-25);
\fill[stB] (1,-24) rectangle (3,-25);
\fill[stA] (3,-24) rectangle (4,-25);
\fill[stG] (4,-24) rectangle (5,-25);
\fill[stB] (5,-24) rectangle (6,-25);
\fill[stA] (6,-24) rectangle (7,-25);
\fill[stB] (7,-24) rectangle (9,-25);
\fill[stA] (9,-24) rectangle (10,-25);
\fill[stB] (10,-24) rectangle (11,-25);
\fill[stA] (12,-24) rectangle (13,-25);
\fill[stG] (13,-24) rectangle (15,-25);
\fill[stA] (0,-25) rectangle (1,-26);
\fill[stB] (1,-25) rectangle (2,-26);
\fill[stG] (3,-25) rectangle (4,-26);
\fill[stG] (5,-25) rectangle (6,-26);
\fill[stB] (6,-25) rectangle (8,-26);
\fill[stB] (9,-25) rectangle (13,-26);
\fill[stG] (13,-25) rectangle (14,-26);
\fill[stB] (14,-25) rectangle (15,-26);
\fill[stA] (0,-26) rectangle (1,-27);
\fill[stB] (1,-26) rectangle (2,-27);
\fill[stG] (2,-26) rectangle (3,-27);
\fill[stB] (3,-26) rectangle (5,-27);
\fill[stA] (5,-26) rectangle (7,-27);
\fill[stB] (8,-26) rectangle (9,-27);
\fill[stB] (10,-26) rectangle (12,-27);
\fill[stG] (12,-26) rectangle (13,-27);
\fill[stB] (13,-26) rectangle (15,-27);
\fill[stA] (0,-27) rectangle (2,-28);
\fill[stB] (2,-27) rectangle (3,-28);
\fill[stA] (3,-27) rectangle (4,-28);
\fill[stB] (6,-27) rectangle (7,-28);
\fill[stG] (7,-27) rectangle (8,-28);
\fill[stA] (8,-27) rectangle (9,-28);
\fill[stB] (9,-27) rectangle (10,-28);
\fill[stG] (11,-27) rectangle (12,-28);
\fill[stB] (12,-27) rectangle (13,-28);
\fill[stA] (13,-27) rectangle (15,-28);
\fill[stF] (0,-28) rectangle (1,-29);
\fill[stA] (2,-28) rectangle (3,-29);
\fill[stB] (3,-28) rectangle (4,-29);
\fill[stG] (5,-28) rectangle (6,-29);
\fill[stA] (7,-28) rectangle (9,-29);
\fill[stB] (9,-28) rectangle (10,-29);
\fill[stG] (10,-28) rectangle (11,-29);
\fill[stA] (11,-28) rectangle (12,-29);
\fill[stG] (12,-28) rectangle (13,-29);
\fill[stF] (14,-28) rectangle (15,-29);

    \draw[white,line width=1.6pt] \staircase{28}{14} -- (15,-14);
    \draw[ink,densely dashed,line width=0.7pt] \staircase{28}{14} -- (15,-14);
    \draw[grid,line width=0.3pt] (0,0) rectangle (15,-29);
    \draw[ink,line width=1.0pt] (0,-28) rectangle (1,-29) (14,-28) rectangle (15,-29);
    \foreach \y/\l in {0/0,14/14,28/28} \node[font=\scriptsize,anchor=east] at (-0.4,-\y-0.5) {$\l$};
    \foreach \x/\l in {1/1,15/15} \node[font=\scriptsize,anchor=north] at (\x-0.5,-29.3) {$\l$};
    \node[font=\small,anchor=south] at (7.5,0.6) {(a) a model, $n=15$};
  \end{scope}
  % (b) the dependency cones of the middle cell and of the right end of the firing row of R_15
  \begin{scope}[shift={(20,-14)}]
    %% generated by figures.py; do not edit
\fill[coneBoth] (14,0) rectangle (15,-1);
\fill[coneBoth] (14,-1) rectangle (15,-2);
\fill[coneBoth] (13,-1) rectangle (14,-2);
\fill[coneBoth] (14,-2) rectangle (15,-3);
\fill[coneBoth] (13,-2) rectangle (14,-3);
\fill[coneBoth] (12,-2) rectangle (13,-3);
\fill[coneBoth] (14,-3) rectangle (15,-4);
\fill[coneBoth] (13,-3) rectangle (14,-4);
\fill[coneBoth] (12,-3) rectangle (13,-4);
\fill[coneBoth] (11,-3) rectangle (12,-4);
\fill[coneBoth] (14,-4) rectangle (15,-5);
\fill[coneBoth] (13,-4) rectangle (14,-5);
\fill[coneBoth] (12,-4) rectangle (13,-5);
\fill[coneBoth] (11,-4) rectangle (12,-5);
\fill[coneBoth] (10,-4) rectangle (11,-5);
\fill[coneBoth] (14,-5) rectangle (15,-6);
\fill[coneBoth] (13,-5) rectangle (14,-6);
\fill[coneBoth] (12,-5) rectangle (13,-6);
\fill[coneBoth] (11,-5) rectangle (12,-6);
\fill[coneBoth] (10,-5) rectangle (11,-6);
\fill[coneBoth] (9,-5) rectangle (10,-6);
\fill[coneBoth] (14,-6) rectangle (15,-7);
\fill[coneBoth] (13,-6) rectangle (14,-7);
\fill[coneBoth] (12,-6) rectangle (13,-7);
\fill[coneBoth] (11,-6) rectangle (12,-7);
\fill[coneBoth] (10,-6) rectangle (11,-7);
\fill[coneBoth] (9,-6) rectangle (10,-7);
\fill[coneBoth] (8,-6) rectangle (9,-7);
\fill[coneBoth] (14,-7) rectangle (15,-8);
\fill[coneBoth] (13,-7) rectangle (14,-8);
\fill[coneBoth] (12,-7) rectangle (13,-8);
\fill[coneBoth] (11,-7) rectangle (12,-8);
\fill[coneBoth] (10,-7) rectangle (11,-8);
\fill[coneBoth] (9,-7) rectangle (10,-8);
\fill[coneBoth] (8,-7) rectangle (9,-8);
\fill[coneBoth] (7,-7) rectangle (8,-8);
\fill[coneRight] (14,-8) rectangle (15,-9);
\fill[coneBoth] (13,-8) rectangle (14,-9);
\fill[coneBoth] (12,-8) rectangle (13,-9);
\fill[coneBoth] (11,-8) rectangle (12,-9);
\fill[coneBoth] (10,-8) rectangle (11,-9);
\fill[coneBoth] (9,-8) rectangle (10,-9);
\fill[coneBoth] (8,-8) rectangle (9,-9);
\fill[coneMid] (7,-8) rectangle (8,-9);
\fill[coneMid] (6,-8) rectangle (7,-9);
\fill[coneRight] (14,-9) rectangle (15,-10);
\fill[coneRight] (13,-9) rectangle (14,-10);
\fill[coneBoth] (12,-9) rectangle (13,-10);
\fill[coneBoth] (11,-9) rectangle (12,-10);
\fill[coneBoth] (10,-9) rectangle (11,-10);
\fill[coneBoth] (9,-9) rectangle (10,-10);
\fill[coneMid] (8,-9) rectangle (9,-10);
\fill[coneMid] (7,-9) rectangle (8,-10);
\fill[coneMid] (6,-9) rectangle (7,-10);
\fill[coneMid] (5,-9) rectangle (6,-10);
\fill[coneRight] (14,-10) rectangle (15,-11);
\fill[coneRight] (13,-10) rectangle (14,-11);
\fill[coneRight] (12,-10) rectangle (13,-11);
\fill[coneBoth] (11,-10) rectangle (12,-11);
\fill[coneBoth] (10,-10) rectangle (11,-11);
\fill[coneMid] (9,-10) rectangle (10,-11);
\fill[coneMid] (8,-10) rectangle (9,-11);
\fill[coneMid] (7,-10) rectangle (8,-11);
\fill[coneMid] (6,-10) rectangle (7,-11);
\fill[coneMid] (5,-10) rectangle (6,-11);
\fill[coneMid] (4,-10) rectangle (5,-11);
\fill[coneRight] (14,-11) rectangle (15,-12);
\fill[coneRight] (13,-11) rectangle (14,-12);
\fill[coneRight] (12,-11) rectangle (13,-12);
\fill[coneRight] (11,-11) rectangle (12,-12);
\fill[coneMid] (10,-11) rectangle (11,-12);
\fill[coneMid] (9,-11) rectangle (10,-12);
\fill[coneMid] (8,-11) rectangle (9,-12);
\fill[coneMid] (7,-11) rectangle (8,-12);
\fill[coneMid] (6,-11) rectangle (7,-12);
\fill[coneMid] (5,-11) rectangle (6,-12);
\fill[coneMid] (4,-11) rectangle (5,-12);
\fill[coneNone] (3,-11) rectangle (4,-12);
\fill[coneRight] (14,-12) rectangle (15,-13);
\fill[coneRight] (13,-12) rectangle (14,-13);
\fill[coneRight] (12,-12) rectangle (13,-13);
\fill[coneNone] (11,-12) rectangle (12,-13);
\fill[coneNone] (10,-12) rectangle (11,-13);
\fill[coneMid] (9,-12) rectangle (10,-13);
\fill[coneMid] (8,-12) rectangle (9,-13);
\fill[coneMid] (7,-12) rectangle (8,-13);
\fill[coneMid] (6,-12) rectangle (7,-13);
\fill[coneMid] (5,-12) rectangle (6,-13);
\fill[coneNone] (4,-12) rectangle (5,-13);
\fill[coneNone] (3,-12) rectangle (4,-13);
\fill[coneNone] (2,-12) rectangle (3,-13);
\fill[coneRight] (14,-13) rectangle (15,-14);
\fill[coneRight] (13,-13) rectangle (14,-14);
\fill[coneNone] (12,-13) rectangle (13,-14);
\fill[coneNone] (11,-13) rectangle (12,-14);
\fill[coneNone] (10,-13) rectangle (11,-14);
\fill[coneNone] (9,-13) rectangle (10,-14);
\fill[coneMid] (8,-13) rectangle (9,-14);
\fill[coneMid] (7,-13) rectangle (8,-14);
\fill[coneMid] (6,-13) rectangle (7,-14);
\fill[coneNone] (5,-13) rectangle (6,-14);
\fill[coneNone] (4,-13) rectangle (5,-14);
\fill[coneNone] (3,-13) rectangle (4,-14);
\fill[coneNone] (2,-13) rectangle (3,-14);
\fill[coneNone] (1,-13) rectangle (2,-14);
\fill[coneRight] (14,-14) rectangle (15,-15);
\fill[coneNone] (13,-14) rectangle (14,-15);
\fill[coneNone] (12,-14) rectangle (13,-15);
\fill[coneNone] (11,-14) rectangle (12,-15);
\fill[coneNone] (10,-14) rectangle (11,-15);
\fill[coneNone] (9,-14) rectangle (10,-15);
\fill[coneNone] (8,-14) rectangle (9,-15);
\fill[coneMid] (7,-14) rectangle (8,-15);
\fill[coneNone] (6,-14) rectangle (7,-15);
\fill[coneNone] (5,-14) rectangle (6,-15);
\fill[coneNone] (4,-14) rectangle (5,-15);
\fill[coneNone] (3,-14) rectangle (4,-15);
\fill[coneNone] (2,-14) rectangle (3,-15);
\fill[coneNone] (1,-14) rectangle (2,-15);
\fill[coneNone] (0,-14) rectangle (1,-15);

    \draw[ink,line width=0.5pt] (15,0) -- (14,0) \foreach \k in {1,...,14} { -- (15-\k,-\k) -- (14-\k,-\k) }
      -- (0,-15) -- (15,-15) -- cycle;
    \node[star,star points=5,fill=white,draw=ink,line width=0.4pt,inner sep=0pt,minimum size=0.2cm] at (7.5,-14.5) {};
    \node[star,star points=5,fill=white,draw=ink,line width=0.4pt,inner sep=0pt,minimum size=0.2cm] at (14.5,-14.5) {};
    \node[font=\small,anchor=south west] at (-1,15.6) {(b) two cones in $R_{15}$};
  \end{scope}
  % legend
  \node[anchor=north west,font=\scriptsize,align=left] at (39,-6) {%
    in (a):\ \swatch{white}\,$\Lq$\ \swatch{stG}\,$\Gq$\ \swatch{stA}\,$\Aq$\ \swatch{stB}\,$\Bq$\
    \swatch{stF}\,$\Fq$\\[0.4ex]
    \phantom{in (a):}\ dashed: the boundary of $R_{15}$\\[1.4ex]
    in (b), the backward dependency cones of the cells $\bigstar$:\\[0.4ex]
    \styleswatch{coneBoth}\ in both cones\\[0.3ex]
    \styleswatch{coneMid}\ only in the cone of the middle cell\\[0.3ex]
    \styleswatch{coneRight}\ only in the cone of the right end\\[0.3ex]
    \styleswatch{coneNone}\ in neither\\[0.6ex]
    middle cell: $76$ of the $120$ cells; right end: $64$};
\end{tikzpicture}

\caption{(a)~The line of length $15$ under the second model above, which
contains the germ of $\delta_{14}^{\mathrm d}$ and satisfies (G1)--(G3) with
$T=7$ for all $n\ge9$ and (G5) for $9\le n\le20$. No cell is in state $\Fq$
before time $2n-2=28$, and at time $28$ only the end cells $1$ and $15$
(outlined) are. (b)~The backward dependency cones, inside $R_{15}$, of the
middle cell ($\kappa=7$) and of the right end ($\kappa=0$) of the firing row
(both marked $\bigstar$), drawn level with (a). The first covers $76$ of the
$120$ cells of $R_{15}$, the second $64$; condition (G5) constrains only the
second.}
\label{fig:model}
\end{figure}

We close with four questions.
\begin{enumerate}[(1)]
\item Does a five-state minimal-time solution exist?
\item Is there a condition on the middle of the reflected triangles that is as
economical as \cref{lem:G}?
\item Can the bound $(k-1)^{2s}\ell$ of \cref{lem:left} be made polynomial in
$s$?
\item Do Balzer's conditions exclude six states?
\end{enumerate}
